\def\stat{gudkov}








\def\tit{МЕТОД ПАРАЛЛЕЛЬНЫХ ЦЕПЕЙ ДЛЯ РАСПОЗНАВАНИЯ 


ИЗОБРАЖЕНИЙ ОТПЕЧАТКОВ ПАЛЬЦЕВ}





\def\titkol{Метод параллельных цепей для распознавания 


изображений отпечатков пальцев} 





\def\autkol{В.\,Ю. Гудков}





\def\aut{В.\,Ю. Гудков$^1$}





\titel{\tit}{\aut}{\autkol}{\titkol}











%{\renewcommand{\thefootnote}{\fnsymbol{footnote}}


%\footnotetext[1]


%{Работа  выполнена при поддержке РФФИ (проект 11-07-00225).}}





\renewcommand{\thefootnote}{\arabic{footnote}}


\footnotetext[1]{Челябинский государственный университет, diana@sonda.ru} 








      \thispagestyle{headings}


      


                  \label{st\stat}


                  


                  \vspace*{-9pt}








\Abst{Предложен новый метод обработки изображений отпечатков 


пальцев, основанный на совместном прослеживании направлений 


папиллярных линий в двух каналах~--- света и тени (вместо прослеживания 


одного канала). В~областях с хорошим качеством изображения направления 


потока в обоих каналах прослеживаются достаточно четко. Но в областях 


плохого качества направления могут прослеживаться заметно хуже. После 


сглаживания изображения и построения светотеней в каждом сегменте изображения формируются 


цепочки точек минимальной и максимальной величины. На основе результатов 


прослеживания формируются гипотезы о направлении потока папиллярных 


линий по каждому каналу, после чего принимается общее решение о 


направлении потока. Прослеживание цепочек в двух каналах поз\-во\-ляет 


существенно увеличить точность измерения направлений и плотностей 


папиллярных линий в областях с низким качеством изображения. 


Предложенный метод устойчив к изменению ширины линий и просветов, 


разрывам линий и другим дефектам изображения. Результатом обработки 


изображения является поле направлений потока папиллярных линий с 


матрицей достоверности. Матрица достоверности рассчитывается по степени 


согласованности направлений в каналах света и тени. Такое уточнение 


направления потока позволяет существенно увеличить точность 


распознавания отпечатков пальцев за счет более качественного выделения 


информативных признаков: окончаний и разветвлений папиллярных линий.}





\vspace*{-2pt}





\KW{отпечаток пальца; поле направлений; светотени; параллельные цепи}





\vspace*{-2pt}





\DOI{10.14357\!/08696527130203}





\vspace*{-12pt}





\section{Введение}


  


  В компьютеризированных системах (КС) идентификацию изображений, как 


правило, выполняют после обработки изображений~[1,~2]. При этом ошибки 


обработки напрямую влияют на ошибки идентификации, которые желательно 


минимизировать~[3,~4].


  





  


  При распознавании дактилоскопических изображений (ДИ) выделяют 


признаки, формирующие структуру самого ДИ. К~таким основным признакам, 


согласно~[5], относят: поле направлений линий, поле плотности линий и поле 


качества линий в виде матриц, размерность которых определяется 


сегментацией\linebreak\vspace*{-12pt}





\pagebreak





\





\vspace*{-12pt}





  \begin{floatingfigure}{66mm} %fig1


      \vspace*{1pt}


 \begin{center}


 \mbox{%


 \epsfxsize=60.795mm


 \epsfbox{gud-1.eps}


 }


 \end{center}


 \vspace*{-6pt}


  \Caption{Изображение отпечатка пальца}


  \vspace*{14pt}


  \end{floatingfigure}





\noindent


 ДИ. В ячейку каждой такой мат\-рицы записывают локально 


вычис\-ленное значение признака (яркость, контраст\-ность, градиент)~[6]. Это 


случайная функция двух аргументов (координат), значение которой неизвестно 


до исхода эксперимента. Очевидно, что случайные поля~[7], генерируемые 


обработкой, определяются распределением яркостей точек ДИ.


  


  На рис.~1 представлено ДИ с дефектами, увеличивающими ошибки 


идентификации~[5]. Это вынуждает разработчиков КС применять новые 


методы распознавания~\cite{2-gk, 3-gk}.





\section{Постановка задачи}





  Изображение как множество действительных чисел формируют в виде 


$F=\{f(x,y)\vert (x,y)\hm\in X\times Y\}$ в прямоугольной области~$G$ 


мощностью $\vert G\vert \hm= x_0y_0$, где $X\hm=0, \ldots ,x_0-1$ и $Y\hm=0, \ldots ,y_0-1$. 


Структурно обработку изображения представляют в виде 


пирамиды~$\mathfrak{R}$ слоев (препаратов изображения) из 


взаимосвязанных иерархий~\cite{4-gk}. Сегментация $l$-го слоя $k$-й 


иерархии $F_k^{(l)}$ разбивает слой на $x_hy_h$ непересекающихся 


квадратных сегментов $S_{hk}^{(l)}(x,y)$ с длиной стороны $2^{h-k}$ и 


вершинами $(x,y)\hm\in X_h\times Y_h$, где $k\hm<h$ и $h$~--- номер 


иерархии; $X_h\hm=0, \ldots ,x_h-1$ и $Y_h\hm=0, \ldots ,y_h-1$. Доступ к каждой точке 


сегмента $S_{hk}(x,y)$ записывают в координатах $(u,v)\hm\in 


\overline{X}_{hk}\times \overline{Y}_{hk}$:


  \begin{equation}


  \left.


  \begin{array}{rl}


  \overline{X}_{hk} &=\{ u+x\cdot 2^{h-k}\vert x\in X_h\wedge u\in 0, \ldots ,2^{h-k}-


1\}\,;\\[9pt]


  \overline{Y}_{hk} &= \{ v+y \cdot 2^{h-k}\vert y\in Y_h\wedge v\in 0, \ldots ,2^{h-k}-1\,,


  \end{array}


  \right\}


  \label{e1-gk}


  \end{equation}


а центры сегментов $(u,v)\hm\in \hat{X}_h\times \hat{Y}_h$ находят в виде: 


\begin{equation}


\left.


\begin{array}{rl}


\hat{X}_h & = \{2^{h-1}+x\cdot 2^h\vert x\in X_h\}\,;\\[9pt]


\hat{Y}_h &= \{ 2^{h-1}+y\cdot  2^h\vert y\in Y_h\}\,.


\end{array}


\right\}


\label{e2-gk}


\end{equation}


  


  Для формализации методов классификационного анализа (КА) применяют 


прямолинейные щелевые $A_h(x,y,\alpha,w)$ и точечные $A_h^o(x,y,\alpha,w)$ 


апертуры. Их как множества в виде элементов $(u,v,\beta)$ определяют по 


формулам:





\noindent


  \begin{equation}


  \left.


  \begin{array}{rl}


  A_h(x,y,\alpha,w) &= \{(u,v,\beta)=(x+w_x,\ y+w_y,\ \beta)\vert w\in 


Z_w\}\,;\\[9pt]


  A_h^o(x,y,\alpha,w) &= \{(u,v,\beta)=(x+w_x,\ y+w_y,\ \beta)\vert w\in 


Z^o_w\}\,,


  \end{array}


  \right\}


  \label{e3-gk}


  \end{equation}


где $w_x=w\cos\alpha$ и $w_y\hm=w\sin\alpha$; $(x,y)\hm\in X_h\times Y_h$~---


центр апертуры; $(u,v)\hm\in X_h\times Y_h$~--- точка апертуры; $w$~--- 


размер апертуры; множества $Z_w\hm=1, \ldots ,w$ и $Z_w^o\hm=\{w\}$; $\alpha$~--- 


угол направления апертуры. Угол, определяющий направление из центра 


$(x,y)$ в точку $(u,v)$ апертуры, находят в виде:


$$


\beta = \arctg \left( \fr{v-y}{u-x}\right) +\pi n$$


при $n\in 0, \ldots ,1$.


  


  Алгоритм движения по точкам ДИ задается отношением:


  \begin{equation}


  R=\{(x_d,y_d)\vert d\in 0, \ldots ,7\}


  \label{e4-gk}


  \end{equation}


полного и строгого порядка по направлению движения~$d$ в 8-свя\-зан\-ной 


области, где $(1,\,0)$ есть элемент для $d\hm=0$.


  


  Отношение $R$ определяет функцию перехода $T_h^d(x,y)$ в $h$-й 


иерархии в направлении~$d$. Функция описывает смещение как отдельной 


точки с координатами $(x,y)\hm\in X_h\times Y_h$, так и синхронное смещение 


множества точек в виде:


  \begin{equation}


  T_h^d\left( \left\{ (x,y)\right\}\right) =\left\{ (a,b)=(x+x_d,\ y+y_d)\vert (a,b)\in 


X_h\times Y_h\right\}\,.


  \label{e5-gk}


  \end{equation}


  


  Задача заключается в построении помехоустойчивого метода измерения 


полей направлений ДИ с опорой на формальный 


  аппарат~(\ref{e1-gk})--(\ref{e5-gk}).


  


  Измерение полей направлений является частью общей задачи обработки ДИ, 


представляемой в виде последовательности процедур и функций, 


сгруппированных поэтапно, очередность выполнения которых определяется 


при разработке КС и реже~--- при обучении. Это увязывает слои 


пирамиды~$\mathfrak{R}$ между собой запрограммированной 


последовательностью отображения данных.


  


\section{Метод параллельный цепей}





  На основе формализма~(\ref{e1-gk})--(\ref{e5-gk}) рассматривается 


коррекция изображения, его сглаживание, построение светотеней и измерение 


полей направлений. Собственно метод параллельных цепей применяется на 


этапе измерения полей направлений как набора случайных величин.


  


\subsection{Коррекция изображения} %3.1.





  Обычной процедурой КА является коррекция исходного изображения 


$F_0^{(0)}\hm= \left\{ f_0^{(0)}(x,y)\right\} \hm= \{f(x,y)\} \hm= F$ для 


обеспечения полного динамического диапазона яркостей:


  \begin{equation}


  F_0^{(1)} = \left\{ f_0^{(1)} (x,y)\right\} =\left\{ \fr{(f_0^{(0)}(x,y)-f_{\min}) 


(2^b-1)}{f_{\max}-f_{\min}}\right\}\,,


  \label{e6-gk}


  \end{equation}


где $f_{\max}$ и $f_{\min}$~--- наибольшее и наименьшее значение яркости 


изображения; $b$~--- глубина изображения~[1]. Расчет $f_{\min}$ и $f_{\max}$ 


можно улучшить, опираясь на гистограмму яркостей или модулей градиента и 


задавая наименьший и наибольший процентили, по которым определяют 


$f_{\min}$ и $f_{\max}$~\cite{3-gk}. В~последнем случае эти параметры 


функции~(\ref{e6-gk}) оценивают на этапе обучения программного объекта.


  


\subsection{Сглаживание изображения} %3.2.


  


  Данная процедура выполняется по формуле двумерной дискретной свертки


  \begin{equation}


  F_0^{(2)} =\left\{ f_0^{(2)}(x,y)\right\} = \left\{ \mathbf{H}** 


f_0^{(1)}(x,y)\right\}


  \label{e7-gk}


  \end{equation}


с ядром вида


$$


\mathbf{H} =\left[ h(i,j)\right] = \begin{bmatrix}


1&2&1\\


2&4&2\\


1&2&1


\end{bmatrix}\,.


$$


Такая фильтрация применяется для расфокусировки изображения, 


преимущественно краев линий~\cite{4-gk}. Это улучшает результаты КА, особенно для 


изображений, снятых с ксерокопий и содержащих обедненный состав серого 


цвета. Утверждать, что сглаживание исходного изображения может увеличить 


ошибки распознавания изображения при его обработке, не следует, так как слои 


$F_0^{(0)}$, $F_0^{(1)}$ и~$F_0^{(2)}$ одновременно хранятся в памяти 


машины, а функции КА имеют доступ к любому слою 


пирамиды~$\mathfrak{R}$.





\begin{figure}[b] %fig2


 \vspace*{1pt}


 \begin{center}


 \mbox{%


 \epsfxsize=127mm %.346mm


 \epsfbox{gud-2.eps}


 }


 \end{center}


 \vspace*{-6pt}


\Caption{Четыре слоя светотеней: (\textit{а})~$F_0^{(0+3)}$; 


(\textit{б})~$F_0^{(1+3)}$; 


(\textit{в})~$F_0^{(2+3)}$;


(\textit{г})~$F_0^{(3+3)}$}


\end{figure}





\subsection{ Построение светотеней} %3.3.





  Сглаженный слой $F_0^{(2)}$ по формуле~(\ref{e7-gk}) служит основой для 


формирования слоев светотеней $F_0^{(d+3)} \hm= \left\{ 


f_0^{(d+3)}(x,y)\right\}$ по формуле двумерной дискретной свертки


  \begin{equation}


  F_0^{(d+3)} =\mathbf{H}^{d}** F_0^{(2)}\,,


  \label{e8-gk}


  \end{equation}


где $d\in D=0, \ldots ,3$~--- четыре направления засветки изображения через 


45$^\circ$; $\mathbf{H}^d$~---  маски Собела~\cite{4-gk} в виде:


$$


\mathbf{H}^0=\begin{bmatrix}


-1&0&1\\ -2&0&2\\ -1&0&1\end{bmatrix}\,;


\enskip \mathbf{H}^1 =\begin{bmatrix}


0&1&2\\ -1 & 0& 1\\ -2 & -1& 0\end{bmatrix}\,;\enskip \ldots ;\enskip


\mathbf{H}^3=\begin{bmatrix}


2&1&0\\ 1&0&-1\\ 0&-1&-2 \end{bmatrix}\,.


$$


Четыре слоя светотеней, вычисленные для ДИ, показаны на рис.~2. Нулевая 


реакция окрашена серым цветом, положительные значения светлее, а 


отрицательные~--- темнее.











\subsection{Измерение полей направлений} %3.4.





  Направление линии, по сути, близко к ориентации простой окрестности~[3]. 


Поскольку светотени преимущественно формируются границами линий, 


применим метод локального адаптивного параллельного сканирования 


указанных слоев вдоль путей <<тени>> и <<света>> в качестве базовой оценки 


матриц направлений. Выполняемая процедура для направления $d\hm\in 


D\hm=0, \ldots ,3$ реализует отображение


  $$


  \Gamma:\ \left\{ S_h^{(d+3)} \right\} \rightarrow \left\{\left\{ 


\Delta_h\right\},\left\{\Lambda_h\right\}\right\}\,,


  $$


где $S_h^{(d+3)}\hm= \left\{ S_h^{(d+3)}(u,v)\vert (u,v)\hm\in \hat{X}_h\times 


\hat{Y}_h\right\}$~---  множество сегментов с центрами по~(\ref{e2-gk}) слоев 


$F_0^{(d+3)}$; $\Delta_h\hm= [\delta_h(x,y)]$~--- слой~$F_h$ как матрица 


направлений на сегментах изображения~$S_h$ с углами $0\hm\leq 


\delta_h(x,y)\hm<\pi$; $\Lambda_h\hm= [\lambda_h(x,y)]$~--- слой как матрица 


достоверностей~$\Lambda_h$ для матрицы направлений~$\Delta_h$; $h$~--- 


номер иерархии, на которой отображаются случайные поля; $h\hm\in 


H\hm=2, \ldots ,n$. Измерения выполняются в двух каналах, означенных символом 


$k\hm\in \{0,\,1\}$: 0~--- канал <<тени>> и~1~--- канал <<света>>.





  \textbf{Первый этап.} Первый этап фиксирует для каждого базового сегмента 


и заданного направления~$d$ четыре позиции:


  \begin{align*}


  p_q^{lk} &= \left( x_q^{lk},y_q^{lk}\right)\in P^{lk}\,;\\


  p_q^{rk} &= \left( x_q^{rk},y_q^{rk}\right)\in P^{rk}\,,


%  \label{e9-gk}


  \end{align*}


где $q$~--- длина цепи; $k\in \{0,\,1\}$~---  метка канала <<тени>> и <<света>>; 


$l$ и~$r$~---  метки левой и правой позиции; $(x,y)$~---  координаты; $P$~---  


цепь. Цепи формируются в слоях $F_0^{(d+3)}$ как выделенные 


последовательности точек $\{p_i \hm= (x_i,y_i)\}$ в направлении $s\hm\in 


G\hm= \{d+1,d+2,d+3\}$, которое в среднем перпендикулярно направлению 


засветки~$d$. Суммирование выполняют по модулю~8. Эти отсчеты 


соответствуют четырем простым цепям на точках изображения, как на 


вершинах графа. Две простые цепи для <<тени>> и две простые цепи для 


<<света>> $P^{lk}$ и~$P^{rk}$ в процессе движения развиваются независимо:


\begin{equation}


\left.


\begin{array}{rl}


P^{lk} &= \left\{ p_i^{lk} = \left( x_i^{lk},y_i^{lk}\right) \vert i\in 0, \ldots ,q\wedge


p_0^{lk}=p^l\right\}\,;\\[9pt]


P^{rk} &= \left\{ p_i^{rk} = \left( x_i^{rk},y_i^{rk}\right) \vert i\in 0, \ldots ,q\wedge 


p_0^{rk}=p^r\right\}\,,


\end{array}


\right\}


\label{e10-gk}


\end{equation}


где $k\in \{0,\,1\}$~---  метка канала; $q$~---   длина цепи; $p^l\hm=(x^l,y^l)$ и 


$p^r\hm= (x^r,y^r)$~---  две стартовые точки для левых и правых простых 


цепей. Обычно стартовые точки~$p^l$ и~$p^r$ располагаются в соседних 


сегментах. Они для сегментов $S_h^{(d+3)}$ с центрами $\{(u,v)\hm\in  


\hat{X}_h\times \hat{Y}_h\}$ по~(\ref{e2-gk}) при $h\hm=2$ определяются методом 


переноса точки $(u,v)$ в направлении~$d$ и зеркальном ему $\overline{d}\hm= 


d\hm+4$ (суммирование по модулю~8) в апертуре по~(\ref{e3-gk}):


\begin{align*}


\left\{ \left( x^l,y^l\right)\right\} & = \left\{ A_0^o(u,v,45d,t)\vert (u,v)\in 


\hat{X}_h\times \hat{Y}_h\right\}\,;\\


\left\{ \left( x^r,y^r\right)\right\} & = \left\{ A_0^o(u,v,45\overline{d},t)\vert (u,v)\in 


\hat{X}_h\times \hat{Y}_h\right\}\,,


\end{align*}


где $t$~---  расстояние, на которое переносится точка из центра сегмента (в 


реализации $t\hm\in 4, \ldots ,7$). Точка остается в центре сегмента, если при переносе 


ее координаты выходят за границы изображения.





  \begin{figure}[b] %fig3


   \vspace*{1pt}


 \begin{center}


 \mbox{%


 \epsfxsize=112.204mm


 \epsfbox{gud-3.eps}


 }


 \end{center}


 \vspace*{-6pt}


  \Caption{Графы функций переходов $T^d$ для $d\in 0, \ldots ,7$}


  \end{figure}


  


  Независимое развитие четырех простых цепей $P^{lk}$ и $P^{rk}$ 


определяется формулой


  \begin{equation}


  f_0^{(d+3)}(x_i,y_i) =f_0^{(d+3)} \left( T_0^s(x_{i-1},y_{i-1})\right)\,,


  \label{e11-gk}


  \end{equation}


где $i>0$; $d\in D=0, \ldots ,3$~---  направление засветки; $T_0^s(\cdot)$~---  функция 


перехода по~(\ref{e5-gk}); $s$~--- направление для функции перехода как 


аргумент выражения


$$


s=\begin{cases}


\mathrm{arg}\,\max\limits_G \left( f_0^{(d+3)}(x_i,y_i)\right) &\ \mbox{при}\ 


k=0\,;\\


\mathrm{arg}\,\min\limits_G \left( f_0^{(d+3)}(x_i,y_i)\right) &\ \mbox{при}\ 


k=1\,.


\end{cases}


$$


  


  Итак, на каждом шаге развития цепи путь исследуется. Графы функций 


переходов показаны на рис.~3.


  


  Этап фиксирует в каждом сегменте для текущего направления~$d$ четыре 


позиции, точность определения которых существенна. Поэтому величина~$q$ 


выбирается оптимальной по критерию: рост $q$ приводит к увеличению 


времени обработки, уменьшение~$q$ снижает точность определения 


направлений даже на плавных линиях из-за наличия на ДИ пор, складок, 


залипаний, рыхлого фона старой бумаги и~т.\,д. Для диагональных 


направлений выбирают целую часть величины  $q^\prime\hm= q\cdot 0{,}707$.


  





  


  Дополнительная оптимизация выполняется шевелением точек~$p^l$ и~$p^r$ 


в виде четырех шагов: шаг в направлении~$d$, два шага в 


направлении~$\overline{d}$ и еще шаг в направлении~$d$ по~(\ref{e11-gk}). 


При этом точки~$p^l$ и~$p^r$ расщепляются и образуются не две, а четыре 


стартовые точки для развития четырех независимых простых цепей 


аналогично~(\ref{e10-gk}).








  \textbf{Второй этап.} На сегментах $S_h$ измеряются параметры 


направлений. Из цепей по~(\ref{e10-gk}) зададим два множества $P^k\hm= 


\left\{ p^{lk},p^{rk}\right\} \hm= \left\{ p_q^{lk},p_q^{rk}\right\}$, $k\hm\in 


\{0,\,1\}$. Каждое множество содержит две точки для <<тени>> либо две точки 


для <<света>>. Расчет сводится к выделению в слоях $F_0^{(d+3)}$ 


по~(\ref{e8-gk}) отсчетов, отвечающих последовательностям точек $\{p_i\hm= 


(x_i,y_i)\}$ в направлении $s\hm\in \overline{G}\hm= \{d+5,d+6,d+7\}$. 


Множество~$\overline{G}$ по направлениям зеркально множеству~$G$, а 


суммирование выполняют по модулю~8. Выбор направлений из~$\overline{G}$ 


позволяет прослеживать светотени в направлении, обратном предшествующему 


этапу. Выделенные точки формируют две простые параллельные цепи для 


<<тени>> и две простые параллельные цепи для <<света>>, причем 


параллельные цепи определяются одинаковой функцией перехода на каждом 


шаге:


  \begin{equation}


  \left.


  \begin{array}{rl}


  \overline{P}^{lk} &= \left\{ p_i^{lk}=\left( x_i^{lk},y_i^{lk}\right)\vert i\in 


q, \ldots ,3q\right\}\,;\\[9pt]


\overline{P}^{rk} &= \left\{ p_i^{rk}=\left( x_i^{rk},y_i^{rk}\right)\vert i\in 


q, \ldots ,3q\right\}\,,


\end{array}


\right\}


\label{e12-gk}


\end{equation}


где $k\in \{0,\,1\}$; $3q-q$~---  длина цепи; $l$ и $r$~---  метки левой и правой 


цепи. Синтез на решетке двух копий простых параллельных цепей 


определяется условием:


\begin{equation}


f_0^{(d+3)} (x_i,y_i) =f_0^{(d+3)}\left( T_0^s(x_{i-1},y_{i-1})\right)\,,


\label{e13-gk}


\end{equation}


где $T_0^s(\cdot)$~---   функция перехода по~(\ref{e5-gk}); начала цепей 


$(x_q,y_q)\hm\in P^k$; $s$~---   направление движения для функции перехода 


как аргумент выражения


$$


s=\mathrm{arg}\,\Psi^k_{\overline{G}} \left( f_0^{(d+3)} \left( x_i^l,y_i^l\right), 


f_0^{(d+3)} \left( x_i^r,y_i^r\right)\right)


$$


с функцией исследования пути на каждом шаге развития цепей:


$$


\Psi^k_{\overline{G}}(a,b) = \begin{cases}


\min\limits_{s\in\overline{G}}\, \max(a,b) &\ \mbox{при}\ k=0\,;\\


\max\limits_{s\in\overline{G}}\, \min(a,b) &\ \mbox{при}\ k=1\,.


\end{cases}


$$








  Для определения направлений <<тени>> и <<света>> на сегменте 


$S_h^{(d+3)}(u,v)$ при рассматриваемых засветках изображения используются 


пары точек $(p_q^{lk},p^{lk}_{3q})$ или $(p_q^{rk},p^{rk}_{3q})$ нулевой 


иерархии. Связанные с ними направления $\delta_n^{(dk)}(x,y)\hm\in F_2$ для 


заданного направления~$d$ определяются по формуле для левой (или правой) 


цепи:





\vspace*{4pt}





\noindent


  \begin{equation*}


  \delta_h^{(dk)}(x,y) =\left( \arctg \left( \fr{y_q^{lk}-y_{3q}^{lk}} {x_q^{lk}-


x_{3q}^{lk}}\right) +\pi n\right) \mathrm{mod}\, \pi \ \mbox{при}\ n\in 0, \ldots ,1\,.


%  \label{e14-gk}


  \end{equation*}


  


  \vspace*{-4pt}


  


  Действительно, выбор точек для расчета направлений цепей 


$\overline{P}^{lk}$ или $\overline{P}^{rk}$ не имеет значения, так как цепи 


канала~$k$ параллельны. Показателем достоверности такого решения могут 


служить величины  $\lambda_h^{(dk)}(x,y)\hm\in F_h$, собираемые в слоях 


данных:





\noindent


  \begin{equation}


  \lambda_h^{(dk)}(x,y) =\fr{1}{q}\left\vert \sum\limits_{i=q}^{3q} \Psi^k \left( 


f_0^{(d+3)} \left( x_i^{lk},y_i^{lk}\right), f_0^{(d+3)} \left( 


x_i^{rk},y_i^{rk}\right)\right)\right\vert\,,


  \label{e15-gk}


  \end{equation}


где точки $f_0^{(d+3)}(x_i,y_i)$ выбираются по~(\ref{e13-gk}); в $k$-м канале 


функция выбора точек из параллельных цепей находится в виде:





\vspace*{3pt}





\noindent


$$


\Psi^k(a,b) =\begin{cases}


\max (a,b) &\ \mbox{при}\ k=0\,;\\


\min (a,b) &\ \mbox{при}\ k=1\,;


\end{cases}


$$





\vspace*{-3pt}





\noindent


$l$ и $r$~---  метки левой и правой цепей. Достоверность~$\lambda_h$ можно 


рассматривать как вероятность, отображенную на шкалу 0--2$^b$ при 


глубине~$b$ изображения.





  Итак, из двух параллельных цепей выбирают одну, а ее начальная и конечная 


вершины определяют направление как угол. Достоверность направления 


оценивают величинами, которые накапливаются в процессе развития цепей. 


Эти величины выбирают из параллельных цепей как значения, наименее 


отклоняющиеся от нуля. Выбор выполняют на каждом шаге развития цепей, 


что напоминает корреляцию цепей, синтезируемых при движении по 


критерию~(\ref{e13-gk}).


  


  Две независимые цепи по~(\ref{e10-gk}), две параллельные цепи (выделены) 


по~(\ref{e12-gk}) и перенос центра сегмента (серые штриховые линии) представлены на рис.~4.


  


  \begin{figure} %fig4


   \vspace*{1pt}


 \begin{center}


 \mbox{%


 \epsfxsize=114.93mm


 \epsfbox{gud-4.eps}


 }


 \end{center}


 \vspace*{-9pt}


  \Caption{Граф цепей для $d=2$, $G=\{3,4,5\}$ в слое <<тени>>}


  \end{figure}


  


  Важность процедуры сглаживания изображения~(\ref{e7-gk}) очевидна. Она 


обеспечивает плавность перепадов яркостей точек. Без этого точность 


измерений ухудшается.


  


  \textbf{Третий этап} необходим для оптимизации элементов матриц на 


базовых сегментах~$S_h$ для заданного $d\hm\in D\hm=0, \ldots ,3$ с помощью 


повторения первого и второго этапов. Заметим, что движение на втором этапе производится 


в секторе, зеркальном сектору первого этапа. И~хотя длина цепи первого этапа 


вдвое короче длины цепи второго этапа, цепи, казалось бы, должны 


центрироваться на сегменте. На самом деле этого не происходит, цепи 


отклоняются, измерения\linebreak\vspace*{-12pt}





\pagebreak





\noindent


 в направлении~$d$ не совпадают с измерениями в 


противоположном направлении $\overline{d}\hm= d\hm+4$. Но 


направления~$d$ и~$\overline{d}$ равноправны, и можно выбрать для 


<<тени>> и <<света>> наилучшее по достоверности измерение.


  


  Оптимизация сводится к вычислению в слоях $F_h$ восьми матриц 


направлений


  \begin{equation}


  \Delta_h^{(dk)} =\left[ \delta_h^{(dk)}(x,y)\right]


  \label{e16-gk}


  \end{equation}


и соответствующих им восьми матриц достоверностей


\begin{equation}


\Lambda_h^{(dk)} =\left[ \lambda_h^{(dk)} (x,y)\right]


\label{e17-gk}


\end{equation}


при $k\in\{0,\,1\}$. При этом в вершину каждого сегмента $S_h^{(d+3)}$ слоя 


$F_h^{(d+3)}$ записываются собственно направления


$$


\delta_h^{(dk)}(x,y) =\delta_h^{(\vartheta (x,y,k)k)}(x,y)


$$


и соответствующие им достоверности


$$


\lambda_h^{(dk)}(x,y) =\lambda_h^{(\vartheta(x,y,k)k)}(x,y)\,,


$$


где направление


\begin{equation}


\vartheta(x,y,k) =\mathrm{arg}\,\max\limits_{\{d,\overline{d}\}} \left( 


\lambda_h^{(dk)}(x,y),\lambda_h^{(\overline{d}k)}(x,y)\right)\,.


\label{e18-gk}


\end{equation}


  


  Направление $\vartheta(x,y,k)\hm\in \left\{d,\overline{d}\right\}$ и 


максимизирует достоверность. Пусть графы цепей такие, как на рис.~4. 


Изменим направление~$d$ на противоположное ему~$\overline{d}$. Несмотря 


на совпадение отсчетов $\{ p^l, p^r\}$, вычисляемых по~(\ref{e10-gk}), графы 


независимых цепей, вероятно, изменятся, вследствие чего и граф параллельных 


цепей тоже изменится. Тогда матрицы направлений и их достоверностей 


оптимизируются. Экспериментально обнаружено, что выбор по~(\ref{e18-gk}) 


улучшает результаты КА.


  


  \textbf{Четвертый этап.} Это процедура перколяции~[2] 


  матриц~(\ref{e16-gk}) и~(\ref{e17-gk}) для направления~$d$ с иерархии 


измерений $h\hm=\min (H)$ на более высокие иерархии 


пирамиды~$\mathfrak{R}$. Процедура сводится к рекурсивному вычислению 


по иерархиям в векторных пространствах восьми матриц направлений


  \begin{equation}


  \Delta_h^{(dk)} =\left[ \delta_h^{(dk)}(x,y)\right] = \left[ \fr{1}{2}\left( \arctg 


\left( \fr{\mathrm{im}_{h-1}^{(dk)}}{\mathrm{re}_{h-1}^{(dk)}}\right) +\pi n\right)\right]


  \label{e19-gk}


  \end{equation}


при $n\in 0,\,1$ и восьми матриц достоверностей


\begin{equation}


\Lambda_h^{(dk)} =\left[ \lambda_h^{(dk)}(x,y)\right] = \left[ \kappa 


\sqrt{(\mathrm{re}_{h-1}^{(dk)})^2+(\mathrm{im}_{h-1}^{(dk)})^2}\,\right]\,,


\label{e20-gk}


\end{equation}


где $k\in 0,\,1$~--- номер канала; иерархия $h\hm\in H\hm=2, \ldots ,n$; направление 


$d\hm\in D\hm= 0, \ldots ,3$. Действительную и мнимую части векторов для сегмента 


$S_{h\,h-1}(x,y)$ определяют по модели сложения векторов с величинами 


модуля $\lambda_{h-1}^{(dk)}(x,y)$ и аргумента $\delta_{h-1}^{(dk)}(x,y)$ в 


виде:


\begin{align*}


\mathrm{re}_{h-1}^{(dk)} &= \sum\limits_{(u,v)\in G}


\lambda_{h-1}^{(dk)} (u,v) \cos \left( 2\delta_{h-1}^{(dk)} (u,v)\right)\,,\\


\mathrm{im}_{h-1}^{(dk)} &= \sum\limits_{(u,v)\in G} \lambda_{h-1}^{(dk)} (u,v) \sin \left( 2\delta_{h-1}^{(dk)} 


(u,v)\right)\,,


\end{align*}


где  $G=\overline{X}_{h\,h-1} \times \overline{Y}_{h\,h-1}$~---  множество 


точек сегмента $S_{h\,h-1}(x,y)$ по~(\ref{e1-gk}); аргумент $\delta_{h-1}^{(dk)} 


(u,v)$ удваивают для расширения полуплоскости~$[0,\,\pi)$ до плоскости 


$[0,\,2\pi)$~\cite{4-gk, 6-gk}; $\kappa$~---  коэффициент нормировки векторов из~$G$ вида


$$


\kappa = \fr{\max\limits_{(u,v)\in G} \lambda_{h-1}^{(dk)}(u,v)} 


{\sum\limits_{(u,v)\in G} \lambda_{h-1}^{(dk)} (u,v)}\,.


$$





  \begin{figure}[b] %fig5


   \vspace*{-6pt}


 \begin{center}


 \mbox{%


 \epsfxsize=126.611mm


 \epsfbox{gud-5.eps}


 }


 \end{center}


 \vspace*{-6pt}


  \Caption{Направления и достоверности канала <<тени>>~(\textit{а}) и канала <<света>>~(\textit{б})}


  \end{figure}


  


  Результат сложения векторов запоминают в вершине сегмента, а данные на 


всех уровнях пирамиды~$\mathfrak{R}$ привязывают к координатам. 


Коэффициент нормировки~$\kappa$ масштабирует достоверность 


$\lambda(x,y)$, которая при $h\hm>2$ как мера когерентности направлений 


достигает максимально возможной величины для идеальной локальной 


ориентации, а для изотропной структуры она равна нулю~\cite{4-gk}.


  


  \textbf{Пятый этап.} Выполняется отбор направлений на сегментах 


$S_h^{(d+3)}$. Формально алгоритм реализуется последовательным выбором в 


каналах направлений из $D\hm=0, \ldots ,3$ для иерархии $H\hm=2, \ldots ,n$ итерационно 


$h\hm=\min(H)\Rightarrow H\hm= H\backslash \{h\}$ по формулам:


  \begin{equation}


  \left.


  \begin{array}{rl}


  \Delta_h &= \left[ \delta_h(x,y)\right] = \left[ 


\delta_h^{(\vartheta(x,y)k)}(x,y)\right]\,; %\label{e21-gk}


\\[9pt]


  \Lambda_h &= \left[ \lambda_h(x,y)\right] = \left[ 


\delta_h^{(\vartheta(x,y)k)}(x,y)\right]\,,


\end{array}


\right\}


\label{e21-gk}


  \end{equation}


где $k\in 0,\,1$ и $\vartheta (x,y,k)\hm\in D$~---  на\-прав\-ле\-ние-по\-бе\-ди\-тель 


на отдельном сегменте иерархии~$h$~---  определяются в виде пары


$$


\left( \vartheta (x,y)k\right) =\mathrm{arg}\,\max\limits_{k\in 0,\,1} 


\max\limits_{d\in D} \lambda_h^{(dk)} (x,y)\,.


$$


  


  Направления покрывают все многообразие ориентаций папиллярных линий 


$0\hm\leq \delta \hm< \pi$ при равноправном выборе направлений~$d$ 


и~$\overline{d}$ по~(\ref{e18-gk}) для простых окрестностей. Матрицы 


$\{\Delta_h\}$, $\{\Lambda_h\}$ содержат наиболее правдоподобные 


на\-прав\-ле\-ния из каждого канала, но необязательно корректные. Здесь для 


каждой иерархии $h\hm\in H$ формируется 18~слоев: $\{ \Delta_h^{(dk)}\}$, 


$\{\Delta_h\}$, $\{\Lambda_h^{(dk)}\}$, $\{\Lambda_h\}$, где $d\hm\in 0, \ldots ,3$; 


$k\hm\in 0,\,1$.








  


  На этом заканчивается измерение полей направлений ДИ. Матрицы 


на\-прав\-ле\-ний и величин их достоверностей показаны на рис.~5 и~6.








  


  Итак, метод измерения полей направлений основан на оценке величин 


отсчетов при параллельном движении в слоях светотеней, а разделение 


результатов КА на два канала и изобилие


 результатов  измерений позволяет\linebreak\vspace*{-12pt}


\begin{floatingfigure}{67mm} %fig6


 \vspace*{1pt}


 \begin{center}


 \mbox{%


 \epsfxsize=61.393mm


 \epsfbox{gud-7.eps}


 }


 \end{center}


 \vspace*{-9pt}


\Caption{Отобранные направления и достоверности}


\end{floatingfigure}  





\noindent


улучшить КА изображения и уменьшить ошибки обработки ДИ и, тем самым, 


уменьшить ошибки идентификации~[5].








  На рис.~6 можно заметить, что, несмотря на дефекты ДИ, обозначенные на 


рис.~1, отобранные на основе максимизации достоверности направления (в 


виде черточек) в этих дефектных областях построены корректно.





\section{Заключение}





  Предложен метод параллельных цепей для распознавания направлений 


линий на ДИ. Он базируется на формализме~(\ref{e1-gk})--(\ref{e5-gk}) и 


выполняется на основе следующих новых решений:


  \begin{itemize}


  \item отслеживание слоев светотеней на основе корреляции точек из 


параллельных цепей~(\ref{e13-gk}), зеркалирования направления движения 


по~(\ref{e16-gk})--(\ref{e18-gk}) и вероятностной оценки достоверностей 


направлений по~(\ref{e15-gk});


  \item многоканальное независимое измерение полей направлений и их 


достоверностей согласно уравнениям~(\ref{e12-gk});


  \item перколяцию матриц направлений и их достоверностей на более высокие 


иерархии пирамиды~$\mathfrak{R}$ по~(\ref{e19-gk}) и (\ref{e20-gk}) в рамках 


пирамидальной модели;


  \item отбор наиболее достоверных направлений в каналах и синтез матриц 


более высокого качества отобранных направлений и их достоверностей 


по~(\ref{e21-gk}). %,~(\ref{e22-gk}).


  \end{itemize}


  


  Предложенный метод существенно отличается от метода простых 


окрестностей~\cite{4-gk} как по точности, так и по качеству измерений.


  


  В дальнейшем планируется развитие метода за счет применения трех 


параллельный цепей в сочетании с методом тензорного анализа простых 


окрестностей.


  


{\small\frenchspacing


{%\baselineskip=10.8pt


\addcontentsline{toc}{section}{Литература}


\begin{thebibliography}{9}


\bibitem{1-gk}


\Au{Гонсалес Р., Вудс Р.} Цифровая обработка изображений~/ Пер. с англ. под 


ред. П.\,А.~Чочиа.~--- М.: Техносфера, 2006. 1072~c.


(\Au{Gonzalez~R., Woods~R.} Digital image processing.~--- 3rd ed.~---


Prentice Hall, 2007. 976~p.)


\bibitem{2-gk}


\Au{Потапов А.\,А., Пахомов А.\,А., Никитин С.\,А., Гуляев~Ю.\,В.} Новейшие 


методы обработки изображений.~--- M.: Физматлит, 2008. 496~с.





\bibitem{4-gk}


\Au{Яне Б.} Цифровая обработка изображений~/ Пер. с англ. 


А.\,М.~Измайловой.~--- М.: Техносфера, 2007. 584~с.


(\Au{J$\ddot{\mbox{a}}$hne~B.} Digital image processing.~--- Springer, 2008. 608~p.)





\bibitem{3-gk}


\Au{Визильтер Ю.\,В., Желтов С.\,Ю., Бондаренко~А.\,В. и~др.} Обработка и 


анализ изображений в задачах машинного зрения.~--- М.: Физматкнига, 2010. 


672~с.





\bibitem{5-gk}


\Au{Bolle~R.\,M., Ratha N.\,K.} Automatic fingerprint recognition systems.~--- N.Y.: Springer-Verlag, 2004. 458~p.


\bibitem{6-gk}


\Au{Maltoni D., Maio D., Jain A.\,K., Prabhakar~S.} Handbook of fingerprint 


recognition.~--- L.: Springer-Verlag, 2009. 496~p.


\bibitem{7-gk}


\Au{Королюк В.\,С., Портенко Н.\,И., Скороход А.\,В., Турбин~А.\,Ф.} 


Справочник по теории вероятностей и математической статистике.~--- М.: 


Наука, 1985. 640~с.


 \end{thebibliography}


} }








\vspace*{12pt}





\hrule





\vspace*{2pt}





\hrule





\def\tit{METHOD OF~PARALLEL CIRCUITS FOR FINGERPRINT IMAGE RECOGNITION}





\def\aut{V.\,Yu.~Gudkov}








\titelr{\tit}{\aut}





\vspace*{3pt}





%\noindent


{\centerline{Chelyabinsk State University, diana@sonda.ru}}











\Abste{New method of fingerprint image processing is described. The method is based on 


simultaneous tracking of ridges and valleys directions, instead of\linebreak\vspace*{-12pt}}





\Abstend{tracking ridges. In areas of good 


quality, the directions can be determined confidently both for ridges and valleys. But in areas of poor 


quality, there is a problem with accurate measuring of ridge flow. Given a pixel on filtered 


fingerprint image, the chains of pixels of minimal and maximal intensity are tracked. Based on the 


results of tracking, the hypotheses on ridge flow direction for ridges and valleys are generated. Then, the 


results are fused in order to obtain estimation of ridge flow direction and density. It allows to


significantly improve the accuracy of ridge flow direction and density in the areas of poor quality. 


Finally,  it is used to form the field of flows and its confidential (quality) matrix. More accurate 


measuring of ridge direction improves fingerprint recognition by more accurate feature extraction: 


ridge endings and bifurcations.}





\KWE{fingerprint; field of flows; chiaroscuro; parallel circuits}





\DOI{10.14357\!/08696527130203}





%\Ack


%\noindent











\renewcommand{\bibname}{\protect\rmfamily References}





 {\small\frenchspacing


{%\baselineskip=10.8pt


%\addcontentsline{toc}{section}{Литература}


\begin{thebibliography}{9}


\bibitem{1-gk1}


\Aue{Gonzales, R., and R. Woods.} 2007. \textit{Digital image processing}. 3rd ed.


Prentice Hall.  976~p.





\bibitem{2-gk-1}


\Aue{Potapov, А.\,А., А.\,А. Pakhomov, S.\,А. Nikitin, and Yu.\,V.~Gulyaev}. 2008. 


\textit{Novejshie metody obrabotki izobrazhenij} [\textit{Newest image processing methods}]. 


M.: Fizmatlit. 496~p.





\bibitem{4-gk-1}


\Aue{J$\ddot{\mbox{a}}$hne, B.} 2005. \textit{Digital image processing}. Springer. 608~p.





\bibitem{3-gk-1}


\Aue{Vizilter, Yu.\,V., S.\,Yu. Zheltov, А.\,V. Bondarenko, \textit{et al}.} 2010. \textit{Obrabotka i 


analiz izobrazhenij v zadachah mashinnogo zreniya} [\textit{Processing and analysis of images in 


technical vision}]. М.: Fizmatkniga. 672~p.





\bibitem{5-gk-1}


\Aue{Bolle, R.\,M., and N.\,K. Ratha}. 2004. \textit{Automatic fingerprint recognition systems}. 


N.Y.: Springer-Verlag. 458~p.


\bibitem{6-gk-1}


\Au{Maltoni, D., D. Maio, A.\,K. Jain, and S.~Prabhakar}. 2009. \textit{Handbook of fingerprint 


recognition}. L.: Springer-Verlag. 496~p.





 \label{end\stat}


 


\bibitem{7-gk-1}


\Au{Korolyuk, V.\,S., N.\,I. Portenko, A.\,V. Skorokhod, and А.\,F.~Turbin}. 


1985. \textit{Spravochnik po teorii veroyatnostej i matematicheskoj statistike} 


[\textit{Handbook of probability theory and mathematical statistics}]. М.: 


Nauka. 640~p. 


 \end{thebibliography}


} }





\renewcommand{\bibname}{\protect\rmfamily Литература}


